\documentclass[11pt,a4paper]{article}
\usepackage[utf8]{inputenc}
\usepackage{amsmath,amssymb}
\usepackage{geometry}
\usepackage{hyperref}
\usepackage{listings}
\usepackage{xcolor}
\usepackage{booktabs}
\usepackage{graphicx}
\geometry{margin=2.5cm}
\hypersetup{colorlinks=true,linkcolor=blue,urlcolor=blue}

\lstset{
    basicstyle=\small\ttfamily,
    backgroundcolor=\color{gray!10},
    frame=single,
    breaklines=true,
    language=Python,
}

\title{\textbf{Mirror -- Chiral Toolkit}\\User Manual}
\author{Mingrui Zuo \& Lifeng Ding\\Xi'an Jiaotong-Liverpool University}
\date{Version 0.6 -- June 2026}

\begin{document}
\maketitle
\tableofcontents
\newpage

% =================================================================
\section{Overview}
% =================================================================

\textbf{Mirror} (Chiral Toolkit) is a molecular-structure analysis platform
built with Streamlit.  It provides six functional tabs for the analysis of
porous crystalline frameworks (CIF) and guest molecules (PDB).

\subsection{Technology Stack}

\begin{tabular}{ll}
    \toprule
    Component & Technology \\
    \midrule
    Frontend & Streamlit (Python 3.12) \\
    3D Visualisation & py3Dmol / stmol / NGL Viewer \\
    Crystallography & pymatgen (SpacegroupAnalyzer, Lattice) \\
    Numerical & NumPy, SciPy (ndimage, signal) \\
    Machine Learning & scikit-learn \\
    Plotting & Matplotlib (Agg backend) \\
    File parsing & Custom CIF / PDB / DEF parsers \\
    \bottomrule
\end{tabular}

% =================================================================
\section{Launching the Application}
% =================================================================

Run from the project root:

\begin{lstlisting}[language=bash]
streamlit run app.py
\end{lstlisting}

Or double-click \texttt{run.bat} on Windows.  The app opens at
\url{http://localhost:8501}.

% =================================================================
\section{Tab 1: PDB Enantiomer (Mirror Generator)}
% =================================================================

Generates mirror images of molecules from PDB or RASPA DEF files.

\subsection{Workflow}

\begin{enumerate}
    \item Upload a \texttt{.pdb} or \texttt{.def} file.
    \item View molecular analysis: formula, atom count, mass,
          bounding box, VDW radii, geometric and mass-weighted centroids.
    \item Select a mirror plane: \textbf{YZ (invert x)},
          \textbf{XZ (invert y)}, or \textbf{XY (invert z)}.
    \item The mirrored structure is displayed alongside the original.
    \item Download the mirrored PDB or mirrored DEF.
\end{enumerate}

\subsection{Notes}

\begin{itemize}
    \item Reflection is performed about the geometric centroid.
    \item For DEF files, only the atom-position block is modified;
          force-field parameters are preserved.
\end{itemize}

% =================================================================
\section{Tab 2: CPA Alignment}
% =================================================================

Canonical Principal Axes (CPA) alignment rotates a molecule so that its
principal axes of inertia align with the Cartesian axes.

\subsection{Workflow}

\begin{enumerate}
    \item Upload a \texttt{.pdb} file.
    \item Toggle hydrogen visibility and render style (stick, sphere/CPK,
          line).
    \item Select method: \textbf{mass-weighted} (inertia tensor) or
          \textbf{geometric} (covariance matrix).
    \item Click \textbf{Standardize Structure}.
    \item Four views are shown: Original, Mirrored (YZ reflection),
          CPA-aligned, and MCPA-aligned.
    \item Each view displays principal moments $(I_1,I_2,I_3)$ and the
          ratio $I_1/I_3$.  Degeneracy or near-spherical shape is flagged.
    \item Download the CPA or MCPA PDB.
\end{enumerate}

For the mathematical theory, see the separate \textit{Mathematics and
Algorithms} document.

% =================================================================
\section{Tab 3: CIF Toolkits}
% =================================================================

The central analysis tab for porous crystals.

\subsection{Structure Visualization}

\begin{itemize}
    \item Upload a \texttt{.cif} file.
    \item Adjust supercell dimensions ($n_a,n_b,n_c$).
    \item Toggle cell box and background colour.
    \item Rendered via NGL Viewer.
\end{itemize}

\subsection{Symmetry Analysis}

Runs \texttt{pymatgen.SpacegroupAnalyzer} across multiple tolerances
(\texttt{symprec} = 0.01, 0.05, 0.1, 0.2).  Reports:

\begin{itemize}
    \item Best space-group symbol and number.
    \item Confidence score (fraction of tolerances in agreement).
    \item Sohncke and enantiomorphic flags.
    \item Full scan table.
\end{itemize}

\subsection{Element-based Channel Map}

Detects 1D channels via 2D Voronoi tessellation in the $ab$-plane.

\subsection{Grid-based Pore Detection \& VDW Heatmap}

A detailed method producing:

\begin{itemize}
    \item Distance-field map with pore centre and radius.
    \item Cylindrical projection of wall atoms (element-coloured).
    \item Smooth Gaussian VDW heatmap on the $\theta$--$z$ plane.
    \item FFT-based pore-shape classification: L2, L3, L4, or L6.
\end{itemize}

\subsection{Turbo Mode}

Ultra-fast alternative using Voronoi tessellation + KDTree.  Suitable for
real-time interactive use.

% =================================================================
\section{Tab 4: Accessible Orientations}
% =================================================================

Places a guest molecule inside a host channel and scans all orientations.

\subsection{Workflow}

\begin{enumerate}
    \item Upload host CIF and guest PDB.
    \item Preview four guest orientations: Original, Mirrored,
          CPA-guest, MCPA-guest.
    \item Toggle CPA alignment and select collision method
          (\textbf{voxel} or \textbf{vdw}).
    \item Extract the centralised channel (periodic replication along $c$).
    \item View fit analysis (pore diameter vs. guest box diagonal).
    \item Run channel accessibility (Mollweide map, $\theta$--$\phi$
          heatmap).
    \item Configure scan ($\theta$ and $\phi$ steps, clash scale,
          mirror plane) and click \textbf{Scan Orientations}.
\end{enumerate}

\subsection{Three-Region Partition}

Rolls the guest envelope sphere through the pore cross-section:

\begin{itemize}
    \item \textbf{Green}: envelope fits --- all orientations OK.
    \item \textbf{Amber}: only centroid fits --- limited orientations.
    \item \textbf{Red}: centroid does not fit.
\end{itemize}

\subsection{Pore Explorer}

Interactively positions a crosshair and computes L/D orientation spheres.
Displays 3D sphere, clearance heatmap, and L vs. D comparison.

% =================================================================
\section{Tab 5: Database Browser}
% =================================================================

Browse the CoRE-COF 1242 database.

\begin{itemize}
    \item Search by name or number.
    \item Table with PLD, LCD, Density, Topology.
    \item Select a COF to view metadata.
    \item \textbf{Import \& Run}: loads the CIF, runs symmetry analysis,
          and executes grid-based pore detection with full visual output.
\end{itemize}

% =================================================================
\section{Tab 6: Information}
% =================================================================

Project goals, planned features, and contact details.

% =================================================================
\section{Batch Scripts}
% =================================================================

Scripts in \texttt{modules/scripts/}:

\begin{tabular}{ll}
    \toprule
    Script & Purpose \\
    \midrule
    \texttt{benchmark\_database.py} & Full-database pore detection \\
    \texttt{batch\_pore\_morphology.py} & Batch morphology computation \\
    \texttt{batch\_fitness.py} & Host--guest fitness screening \\
    \texttt{spinol\_ml\_analysis.py} & ML regression on spinol selectivity \\
    \texttt{plot\_morphology\_scatter.py} & CI-coloured scatter plots \\
    \texttt{plot\_topology\_pie.py} & Layer-cake topology/SG/sgType \\
    \texttt{plot\_sgtype\_ddg.py} & sgType vs $\Delta\Delta G$ \\
    \texttt{scaling\_benchmark.py} & Timing-scaling benchmarks \\
    \bottomrule
\end{tabular}

% =================================================================
\section{Output Data}
% =================================================================

Batch results are in \texttt{modules/database/benchmark/}:

\begin{itemize}
    \item \texttt{pore\_morphology.xlsx}: 1099 COFs $\times$ 24 columns
          (pore radius, CI, CF, curvature, Fourier powers, helicity, etc.)
    \item \texttt{benchmark.xlsx}: 1242 COFs with symmetry and dimensions.
    \item PNG figures (topology cake plots, SG/sgType distributions,
          scatter plots, ML diagnostics).
\end{itemize}

\end{document}
